% IJCAR 2026 talk --- 20 minute slot (+5 Q&A).
%
% Flow: three parts. Part 1 is the on-ramp --- SMT solvers, the SMT-LIB
% standard, proof production, Eunoia/Ethos, CPC, LambdaPi (with the
% interop hub), motivation, then eo2lp itself (pipeline, then the demo
% pair as its payoff). Part 2 unpacks the two formalisms and the
% translation; Part 3 is results.
% Design language lives in talk-ijcar-design.tex.
% Voice is label-style: slides carry names, artifacts and numbers;
% the argument is spoken, not printed.
%
% Build:  latexmk -lualatex talk-ijcar.tex
% `handout` collapses overlays to one slide each --- on while we review
% the deck; drop it to get the incremental builds back for presenting.
%
\documentclass[handout, t, lualatex, compress, 11pt]{beamer}
\usetheme{metropolis}
% Our own \divider frames open each section, so suppress metropolis's
% automatic section pages or every part is announced twice.
\metroset{sectionpage=none, numbering=fraction, progressbar=frametitle}

\usepackage{fontspec}   % not all themes load it; \newfontfamily needs it
\usepackage{xcolor}
\usepackage{xcolor-material}
\usepackage{ebproof}
\usepackage{unicode-math-input}
\usepackage{MnSymbol}
\usepackage{stmaryrd}
\usepackage{mathtools}
\usepackage{amsmath}
\usepackage{amssymb}

\usepackage{listings}
\usepackage{array}
\usepackage{booktabs}

% Broken by hand: left to itself the first line fills and the second sits
% at 53%, which is the ragged wrap this deck avoids everywhere else.
\title{Automatically Translating Proof Systems\\for SMT Solvers to the $λΠ$-calculus}
\author{Ciarán Dunne \and Guillaume Burel}
\institute{INRIA, ENS Paris-Saclay \quad ensIIE \quad SAMOVAR}
\date{IJCAR 2026, Lisbon}

\input{talk.macros}

% `unicode-math-input` maps ⤳ to \rightcurvedarrow and ⟦⟧ to
% \lBrack/\rBrack; no loaded math font defines them.
\providecommand{\rightcurvedarrow}{\leadsto}
\providecommand{\lBrack}{\llbracket}
\providecommand{\rBrack}{\rrbracket}

% ---------------------------------------------------------------------
% Glyph discipline. Generated .lp files use 14 non-ASCII codepoints;
% the awkward ones are ⤳ (U+2933) and 𝔹 (U+1D539). listings' `literate`
% does not match these codepoints at all, so the font must contain
% them. Of the installed monospace fonts, DejaVu/FreeMono miss ⤳ and 𝔹,
% JetBrains Mono misses ⤳, ⋅ and ≐. JuliaMono covers all 14.
%
% Dependency: JuliaMono at ~/.local/share/fonts/JuliaMono (v0.63.2)
%   https://github.com/cormullion/juliamono
% ---------------------------------------------------------------------
% Scale=1.0: the talk asks the audience to READ generated LambdaPi, so
% code sits at near-parity with body text rather than at the ~55% it
% lands on with \scriptsize at a smaller scale. Listings never go below
% \footnotesize; anything that then overflows gets trimmed, not shrunk.
\newfontfamily\lpmono{JuliaMono}[Scale=1.0]

% τ is a LambdaPi identifier; \mathsf{τ} renders empty in this deck's
% math font, so set it in the code font. Recoloured in the design file.
\newcommand{\tauop}{\text{\lpmono τ}}

% ---------------------------------------------------------------------
% Design language for the IJCAR talk.
%
% The chrome is FLoC 26's (floc26.org, sampled from the site assets):
% royal blue #1D3C8A, azulejo gold #F6B655, cream #F4F2F0. The deck's
% one idea --- colour encodes SOURCE vs TARGET --- keeps its own green
% for the source language and takes the brand blue for the target:
%
%   green   = Eunoia   (source)   the pre-rebrand teal
%   blue    = LambdaPi (target)   the brand blue itself
%   gold    = the FLoC accent: chrome (progress bar, rules, tile
%             band) and \alert emphasis (darkened for text); never a
%             language, so it can sit next to either
%   ink     = everything else
%   muted   = supporting detail, comments, grammar delimiters
%
% Chrome comes from the site: solid blue frametitle band with cream
% text (their section bands), gold progress bar (their accents), and
% the azulejo tile strip as a footer on every page --- floc-tiles.png,
% cut seamlessly from the site's own strip asset; the white letter
% tiles hidden in it are theirs. The frame number sits in a cream
% square at the band's right end, echoing those letter tiles.
%
% Consequences, enforced below:
%   - the header band and footer tiles are the only blue chrome; body
%     structure (bullets, alerts) stays ink so blue keeps meaning
%     "LambdaPi" in content
%   - no red anywhere; \alert is bold in darkened FLoC gold
%   - listings are two-tone, not the 5-colour smtlib rainbow
%   - no line numbers
%   - title and divider frames sit on FLoC cream (their hero band);
%     content frames stay near-white
% ---------------------------------------------------------------------

\usepackage{tikz}
\usetikzlibrary{arrows.meta, positioning, fit, backgrounds, shadows.blur}

% ---- measure, margins and vertical rhythm ---------------------------
% At 12pt the line measure was ~56 characters, below the comfortable
% 60-75 band --- which is why the type read as oversized and crowded at
% the same time. 11pt with slightly tighter margins gives ~63.
\setbeamersize{text margin left=8mm, text margin right=8mm}

% Body text indented 24pt from the frametitle's left edge looked like a
% misalignment rather than a hierarchy. Pull the lists back under it.
\setlength{\leftmargini}{12pt}
\setlength{\leftmarginii}{12pt}
\setlength{\leftmarginiii}{12pt}

% One spacing unit, derived from the body size, replacing the ad-hoc
% mix of \vspace{1mm}/{2mm}/{3mm} that was scattered through the deck.
\newlength{\gapunit}
\setlength{\gapunit}{0.6\baselineskip}
\newcommand{\gap}{\vspace{\gapunit}}

% Listings sit in the same rhythm instead of carrying their own skips.
\lstset{aboveskip=\gapunit, belowskip=\gapunit}

% Same for booktabs tables: the top rule carries its own air, so a table
% under a caption line needs no \vspace of its own. (Half a unit, not a
% full one --- the rule is already doing part of the separating.)
\setlength{\abovetopsep}{0.5\gapunit}

% metropolis pins the frametitle's leftskip to its own padding length
% (2.2ex), independent of the text margin, so the title hangs ~12pt to
% the left of every body element. Measured on p.21: title at x=23px,
% body at x=48px. That reads as a misalignment rather than a hierarchy.
% Re-issue the template with the leftskip tied to the text margin,
% leaving the padding length alone --- it also drives the vertical
% struts that give the title block its height.
\makeatletter
\defbeamertemplate{frametitle}{aligned}{%
	\nointerlineskip%
	\begin{beamercolorbox}[%
			wd=\paperwidth,%
			sep=0pt,%
			leftskip=\beamer@leftmargin,%
			rightskip=\beamer@rightmargin,%
		]{frametitle}%
		\metropolis@frametitlestrut@start%
		\insertframetitle%
		\nolinebreak%
		\metropolis@frametitlestrut@end%
	\end{beamercolorbox}%
}
\setbeamertemplate{frametitle}[aligned]
\makeatother

% metropolis implements `progressbar=frametitle` by APPENDING the bar to
% whatever the frametitle template is when \metroset runs; re-issuing the
% template above therefore threw it away, and the deck shipped with no
% progress indicator at all. Re-append it here.
\addtobeamertemplate{frametitle}{}{%
	\usebeamertemplate*{progress bar in head/foot}}

% ---- palette --------------------------------------------------------
% FLoC 26 brand colours, sampled from floc26.org assets. The gold is
% 1.8:1 against white --- chrome, never text.
\definecolor{flocblue}{HTML}{1D3C8A}    % site bands, tile blue
\definecolor{flocamber}{HTML}{F6B655}   % tile gold
\definecolor{floccream}{HTML}{F4F2F0}   % hero background
\colorlet{eocol}{MaterialTeal700}     % Eunoia
\colorlet{eobg}{MaterialTeal50}
\colorlet{lpcol}{flocblue}            % LambdaPi
\colorlet{lpbg}{flocblue!6}
\colorlet{ink}{MaterialGrey900}
\colorlet{muted}{MaterialGrey600}
\colorlet{rule}{MaterialGrey400}

% ---- inline code ----------------------------------------------------
% Prose used bare \texttt, which reaches for the theme's Fira Mono at the
% full body size: measured 10.8pt of ink against the 9pt of the JuliaMono
% inside the boxes, so a passing mention of a tool name shouted louder
% than the generated code the audience is asked to read. One family for
% all code, scaled to the body's x-height so it also behaves in
% frametitles, and the source/target colour rule applies inline too.
%
% \text{} rather than a bare font switch: \kw and friends are used inside
% math displays, where \newfontfamily switches are invalid.
% (\lpcode would clash with LuaTeX's protrusion primitive; the tt suffix
% follows \mtt, which the paper's macros already use for the same job.)
% Ligatures off: JuliaMono is a programming font and renders -> as a
% single arrow glyph, but \kw{->} is the literal Eunoia token. (The
% listings font escapes this because listings sets every character
% separately, which breaks the ligature lookup.)
% -calt as well as the ligature keys: JuliaMono builds -> from contextual
% ALTERNATES, which Ligatures=NoContextual (i.e. -clig) does not touch.
\newfontfamily\lpmonoinline{JuliaMono}[Scale=MatchLowercase,
	Ligatures={NoCommon, NoContextual}, RawFeature={-calt}]
\newcommand{\code}[1]{\text{\lpmonoinline #1}}    % tools, flags, output
\newcommand{\eott}[1]{\text{\lpmonoinline\color{eocol}#1}}   % Eunoia
\newcommand{\lptt}[1]{\text{\lpmonoinline\color{lpcol}#1}}   % LambdaPi

% ---- kill the inherited chrome --------------------------------------
% talk.macros sets structure to indigo and recolours per section; make
% it neutral once, here, so no \setbeamercolor is needed per section.
% Body structure stays ink (not brand blue): in content, blue means
% "LambdaPi", and a bullet must not borrow it.
\setbeamercolor*{structure}{fg=ink}
\setbeamercolor*{normal text}{fg=ink}
\setbeamercolor*{section in head/foot}{fg=muted}
\setbeamercolor*{subsection in head/foot}{fg=muted}

% ---- FLoC chrome ----------------------------------------------------
% Header: the site's solid blue band, cream text. The aligned template
% above already spans \paperwidth, so setting bg is all it takes.
\setbeamercolor*{frametitle}{fg=floccream, bg=flocblue}
% Since the header is orientation (section · subsection), not content,
% it cedes its size: \small against metropolis's \large, and the
% band's padding comes down with it. The font alone barely moves the
% band (10.4mm -> 9.9): the struts are padding-dominated, and the
% padding is a fixed length. Trimming it is safe now --- the aligned
% frametitle template above no longer uses it for leftskip, only the
% struts read it. Measured band: 10.4mm at \large/2.2ex, 8.5mm here.
\setbeamerfont{frametitle}{size=\small}
\makeatletter
\setlength{\metropolis@frametitle@padding}{1.5ex}
\makeatother

% Progress bar under the band: gold fill on a pale blue track, thick
% enough to read as an accent rather than a hairline.
\setbeamercolor{progress bar}{fg=flocamber, bg=flocblue!18}
\makeatletter
\setlength{\metropolis@progressinheadfoot@linewidth}{1.4pt}
\setlength{\metropolis@titleseparator@linewidth}{1.4pt}
\makeatother

% Title page: blue title, gold separator rule (metropolis draws it).
\setbeamercolor{title}{fg=flocblue}
\setbeamercolor{title separator}{fg=flocamber}

% ---- frame headers name the position, not the slide -----------------
% Convention (Ciarán, 2026-07-26): the frametitle is the current
% section, qualified by the current subsection when one is set ---
% orientation lives in the header, specifics live on the slide and in
% the spoken script. Every content frame opens
% \begin{frame}{\secheader}; a \subsection{...} before a run of frames
% qualifies them, \subsection*{} returns to the bare section header.
% The subsection sits in a cream-blue blend: readable on the band,
% clearly subordinate to the section. Emptiness is tested by width ---
% \insertsubsectionhead has no reliably \ifx-able empty form.
\newlength{\subsecheadw}
\newcommand{\secheader}{%
	\insertsectionhead
	\settowidth{\subsecheadw}{\insertsubsectionhead}%
	\ifdim\subsecheadw>0pt
		{\color{floccream!70!flocblue}\ ·\ \insertsubsectionhead}%
	\fi}

% Footer: the FLoC azulejo strip, full-bleed at the bottom of every
% frame (plain ones included --- on the site the strip separates
% sections, so dividers keep it). Drawn on the background canvas, so
% it reserves no layout space; floc-tiles.png is baked at exactly
% \paperwidth : \tilebandh (128:4), so width x height is not a stretch.
\newlength{\tilebandh}
\setlength{\tilebandh}{4mm}
\setbeamertemplate{background canvas}{%
	\usebeamercolor{background canvas}%
	\begin{tikzpicture}
		\useasboundingbox (0,0) rectangle (\the\paperwidth,\the\paperheight);
		\ifbeamercolorempty[bg]{background canvas}{}{%
			\fill[color=bg] (0,0) rectangle (\the\paperwidth,\the\paperheight);}
		\node[anchor=south west, inner sep=0pt] at (0,0)
		{\includegraphics[width=\paperwidth, height=\tilebandh]{floc-tiles}};
	\end{tikzpicture}}

% Frame number: a cream square sitting on the strip's right end, blue
% numeral --- the strip hides white letter tiles, this is the deck's
% own. Replaces the metropolis footline (numbering=fraction), whose
% ~11mm of reserved height the band does not need; the text block only
% grows. Plain frames (title, dividers) get no number, as before.
% 0.3pt taller than the band: at some rasterisations the image and the
% node round to different pixel heights and a hairline of pattern shows
% over the tile; the overhang lands cream-on-white, invisible.
\setbeamertemplate{footline}{%
	\begin{beamercolorbox}[wd=\paperwidth, sep=0pt, right]{footline}%
		\tikz\node[fill=floccream, minimum width=\tilebandh,
			minimum height={\dimexpr\tilebandh+0.3pt\relax}, inner sep=0pt,
			font=\tiny, text=flocblue] {\insertframenumber};%
	\end{beamercolorbox}}

% \alert wears the FLoC gold: with Eunoia back on teal no language
% claims it, so the brand accent is free for emphasis (the site uses
% its orange the same way). The tile gold itself is 1.8:1 on white ---
% an accent, not a text colour --- so alerts take the same hue
% darkened to 5.2:1. Weight comes along: gold alone at \small body
% sizes is not enough of a signal.
\definecolor{alertgold}{HTML}{96610D}
\setbeamercolor*{alerted text}{fg=alertgold}
\setbeamerfont{alerted text}{series=\bfseries}

% \mtt marks Eunoia symbols, so it takes the Eunoia colour (was magenta)
% and the code family (was the theme's mono at body size).
\renewcommand{\mtt}[1]{\eott{#1}}

% "Example." led in MaterialGreen900; make it neutral and quiet.
\renewcommand{\exxample}{{\itshape\color{muted}Example.\ }}

% Grammar annotations were MaterialIndigo500 --- indigo now means
% "LambdaPi", so these must not borrow it.
\renewcommand{\mcomment}[1]{{\color{muted}{\text{#1}}}}

% \smt sets language=smtlib explicitly, which resets to the language's
% own 5-colour keyword styles. Pin it to the two-tone style.
\renewcommand{\smt}[1]{\lstinline[style=eo, basicstyle=\lpmonoinline]{#1}}

% ---- semantic text macros -------------------------------------------
\newcommand{\eoterm}[1]{{\color{eocol}#1}}   % a Eunoia thing, in prose
\newcommand{\lpterm}[1]{{\color{lpcol}#1}}   % a LambdaPi thing, in prose
\newcommand{\eoname}{\eoterm{\textbf{Eunoia}}}
\newcommand{\lpname}{\lpterm{\textbf{LambdaPi}}}
\renewcommand{\LambdaPi}{\lpname}
% τ appears inside math displays, so it takes the inline scale, not the
% listings scale --- at Scale=1.0 it sat a size above the surrounding math.
\renewcommand{\tauop}{\lptt{τ}}

% ---- pass/fail marks (the only pictorial marks in the deck) ---------
% Deliberately uncoloured: teal and indigo mean "Eunoia" and
% "LambdaPi", and a status glyph must not borrow either.
\newcommand{\yes}{{\color{ink}$\checkmark$}}
\newcommand{\no}{{\color{muted}$\times$}}

% A claim item: a list entry marked ✓ rather than bulleted.
% Plain \item[\yes] misaligns. beamer \llap's item labels, so the ink's
% RIGHT edge is pinned \labelsep before the text and a label wider than
% the bullet grows leftwards: the checkmark started 6pt left of the
% frametitle, outside the deck's grid entirely. A negative kern after the
% mark pulls the whole label back into the marker column (measured: ink
% from 57px to 68px, against 74px for the bullets and 63px for the title).
% Takes beamer's overlay spec the way \item does: \claim<+-> ...
\newlength{\claimpull}
\setlength{\claimpull}{4pt}
\makeatletter
\newcommand{\claim@mark}{\yes\hspace*{-\claimpull}}
\def\claim@ov<#1>{\item<#1>[\claim@mark]}
\newcommand{\claim}{\@ifnextchar<{\claim@ov}{\item[\claim@mark]}}
\makeatother

% ---- listings: two-tone, not a rainbow ------------------------------
\lstdefinelanguage{lambdapi}{
	morekeywords={symbol, constant, sequential, injective, rule, with,
			require, open, assert, let, in, notation, unif_rule, builtin,
			opaque, protected, private, inductive, TYPE},
	morekeywords=[2]{Set, Proof, Type},
	alsoletter={:,-,>,|,$},
	sensitive=true,
	morecomment=[l]{//},
	morestring=[b]",
}

% Eunoia source listings.
\lstdefinestyle{eo}{
	language=smtlib,
	basicstyle=\lpmono\small\color{ink},
	keywordstyle=\color{eocol}\bfseries,
	keywordstyle=[2]\color{ink},
	keywordstyle=[3]\color{eocol},
	keywordstyle=[4]\color{ink},
	keywordstyle=[5]\color{eocol}\bfseries,
	commentstyle=\color{muted}\itshape,
	stringstyle=\color{muted},
	numbers=none,   % talk.macros bakes numbers=left into `smtlib` itself
	extendedchars=true,
	showstringspaces=false,
}

% Generated LambdaPi listings.
\lstdefinestyle{lp}{
	language=lambdapi,
	basicstyle=\lpmono\small\color{ink},
	keywordstyle=\color{lpcol}\bfseries,
	keywordstyle=[2]\color{lpcol},
	commentstyle=\color{muted}\itshape,
	stringstyle=\color{muted},
	numbers=none,   % talk.macros bakes numbers=left into `smtlib` itself
	extendedchars=true,
	showstringspaces=false,
}

% No line numbers: nothing in the deck ever refers to one, they were the
% only reason the snippet boxes needed 14pt of left padding against 3pt of
% right, and on a one-line snippet a lone "1" is pure noise.

% SMT-LIB input is neither the source nor the target language. The
% pipeline figure already draws .smt2 as a neutral node; its listing has
% to agree, or teal ends up meaning "Eunoia, and also this other thing".
\lstdefinestyle{smt2}{
	style=eo,
	keywordstyle=\color{ink}\bfseries,
	keywordstyle=[2]\color{ink},
	keywordstyle=[3]\color{ink},
	keywordstyle=[4]\color{ink},
	keywordstyle=[5]\color{ink}\bfseries,
}

\lstset{style=eo}

% ---- snippet boxes --------------------------------------------------
% The paper wraps every code fragment in an eobox/lpbox (preamble.tex).
% Adopted here, with one change: the paper's snippet style is a
% six-colour scheme (pink commands, blue attributes, teal builtins,
% purple types), which would reimport the rainbow this deck removed.
% Instead the BOX carries the language identity via its tint and rule,
% and the code inside keeps the two-tone discipline. The container does
% the semantic work, so the type can stay quiet.
\usepackage{tcolorbox}
\tcbuselibrary{skins}

\colorlet{eoboxbg}{eocol!5}
\colorlet{lpboxbg}{lpcol!5}

% listings inserts trivlist glue inside the box; zero it out.
\newcommand{\snipsetup}{\topsep=0pt \partopsep=0pt \parsep=0pt \parskip=0pt
	\lstset{aboveskip=0pt, belowskip=0pt}}

% Padding is symmetric now that the line-number gutter is gone, and the
% skips are a full \gapunit: at 0.5 they measured 5pt against a 13.6pt
% body line, so prose sat closer to the code frames than to itself --- on
% p16 a single line was 5pt from the box above and 5.4pt from the box
% below, captioning both.
%
% The skips are not free: measured to the text block (not the page
% bottom), the fullest frames had ~1-8pt of slack, so p2, p4 and p20 each
% had to give the space back --- by dropping a \gap that the larger skip
% now provides, and on p2 by joining the continuation lines. Anything
% added to those frames from here needs the same treatment.
\tcbset{snippet/.style={enhanced, arc=0pt, boxrule=0.6pt,
			left=6pt, right=6pt, top=1pt, bottom=1pt,
			before skip=\gapunit, after skip=\gapunit}}

\newenvironment{eobox}{%
	\begin{tcolorbox}[snippet, colback=eoboxbg, colframe=eocol!40]\snipsetup
			}{\end{tcolorbox}}

\newenvironment{lpbox}{%
	\begin{tcolorbox}[snippet, colback=lpboxbg, colframe=lpcol!40]\snipsetup
			}{\end{tcolorbox}}

% Neutral container, for code that belongs to neither language.
\newenvironment{plainbox}{%
	\begin{tcolorbox}[snippet, colback=white, colframe=rule]\snipsetup
			}{\end{tcolorbox}}

% Thin rule between fragments sharing one box (as in the paper).
% \hrule, not \rule-in-a-paragraph: a paragraph line carries
% \baselineskip, which opened almost a full line of air ABOVE the
% hairline and none below --- both split listings read top-heavy.
% Spacing follows the booktabs convention (0.4ex above a midrule,
% 0.65ex below): more GLUE below than above, because the line above
% already contributes its mostly-empty descender band while the caps
% below reach the top of their box --- the INK gaps then read equal.
% Scaled up ~2x from booktabs' table-density values and then tuned on
% a 300dpi render: 0.9/1.05ex measures ~2.5mm of ink gap on each side.
\newcommand{\listsep}[1][eocol]{%
	\par\vspace{0.9ex}%
	{\color{#1!30}\hrule height 0.3pt}%
	\vspace{1.05ex}}

% ---- abstract syntax vocabulary -------------------------------------
% Policy: math displays carry ABSTRACT syntax (metavariables, schemas,
% operators); concrete code only ever appears inside a snippet box.
% Grammar terminals are still concrete tokens, marked with \kw/\attr so
% they read as terminals rather than as pasted code.
\newcommand{\kw}[1]{\text{\lpmonoinline\bfseries\color{eocol}#1}}
\newcommand{\attrib}[1]{\eott{:#1}}
\newcommand{\lpkw}[1]{\text{\lpmonoinline\bfseries\color{lpcol}#1}}
\newcommand{\epar}[1]{%
	\wrapp{muted}{\text{\lpmonoinline(}}{\text{\lpmonoinline)}}{#1}}

% ---- displays -------------------------------------------------------
% $$...$$ centres on the slide and takes TeX's own fixed display skips, so
% each display landed on a left edge of its own --- measured 161pt, 258pt
% and 298pt on p14 alone --- while the frametitle and every body element
% share one edge. Displays join that edge and the \gapunit rhythm.
% Inside a list it flushes to the item's text edge, which is what you
% want there. Also keeps abstract syntax out of \[...\], which cannot be
% left-aligned without fleqn.
\newenvironment{disp}
	{\par\vspace{\gapunit}\noindent$\displaystyle}
	{$\par\vspace{\gapunit}}

% ---- section dividers: typographic, not full-bleed -------------------
% Opening a beamer section has no visual effect (metropolis has no
% headline and sectionpage is off) --- it is kept for the PDF outline.
% FLoC-styled: cream ground (the site's hero band), blue title, gold
% rule; the tile strip arrives via the background canvas. The old
% per-part colour argument is gone --- part colours would fight the
% brand scheme.
% #1 part label, #2 title
\newcommand{\divider}[2]{%
	\section{#2}%
	\begingroup
	\setbeamercolor{background canvas}{bg=floccream}%
	\begin{frame}[plain]
		\vfill
		\hspace{0.12\textwidth}%
		\begin{minipage}{0.8\textwidth}
			{\color{muted}\normalsize #1}\\[1.5mm]
			{\color{flocblue}\Huge\textbf{#2}}\\[3mm]
			{\color{flocamber}\rule{0.3\textwidth}{1.4pt}}
		\end{minipage}
		\vfill
	\end{frame}
	\endgroup}

% ---- the pipeline diagram -------------------------------------------
% Columns carry the language (teal = Eunoia, indigo = LambdaPi), so the
% two eo2lp arrows are parallel and it is visibly the same translation
% applied to a signature and to a proof. Rows carry the frequency:
% the signature is translated once, proofs once per problem.
\newcommand{\pipelinefig}{%
	\resizebox{\textwidth}{!}{%
		\begin{tikzpicture}[
			x=1cm, y=1cm,
			box/.style={draw, rounded corners=1.5pt, align=center,
					font=\footnotesize, inner xsep=4pt, inner ysep=3.5pt,
					minimum height=7.5mm, minimum width=21mm},
			eo/.style={box, draw=eocol, fill=eobg, text=eocol},
			lp/.style={box, draw=lpcol, fill=lpbg, text=lpcol},
			neutral/.style={box, draw=rule, text=ink, minimum width=16mm},
			tool/.style={font=\scriptsize\itshape, text=muted,
					inner sep=1.5pt, fill=white},
			side/.style={font=\scriptsize\itshape, text=muted, anchor=east},
			ar/.style={-{Stealth[length=4.5pt]}, draw=muted, semithick},
		]
		% --- signature row (once) ---
		\node[eo] (cpc)  at (3.2, 1.45) {CPC signature};
		\node[lp] (lcpc) at (7.0, 1.45) {CPC encoding};

		% --- proof row (per problem) ---
		\node[neutral] (smt)  at (0.0, 0) {\texttt{.smt2}};
		\node[eo]      (prf)  at (3.2, 0) {proof script};
		\node[lp]      (lprf) at (7.0, 0) {\texttt{.lp} proof};
		\node[neutral] (ok)   at (10.4, 0) {\yes\ checked};

		\draw[ar] (cpc)  -- node[tool, above] {eo2lp} (lcpc);
		\draw[ar] (smt)  -- node[tool, above] {cvc5} (prf);
		\draw[ar] (prf)  -- node[tool, above] {eo2lp} (lprf);
		\draw[ar] (lprf) -- node[tool, above] {lambdapi} (ok);

		% the encoded signature is what the proof is checked against
		\draw[ar, dashed] (lcpc) -- node[tool, right, xshift=1pt]
		{checked against} (lprf);

			\node[side] at (-0.55, 1.45) {once};
			\node[side] at (-0.55, 0)    {per problem};
		\end{tikzpicture}}}

% ---- the SMT-LIB page stack -----------------------------------------
% Real pages of the standard, fanned about a shared bottom-left corner.
% The artefact makes the "this is a 112-page document" point without
% saying it (it had a "v2.6, 104 pages" caption; cut). Static by
% design: the stack is scene-setting, not an argument being built, so
% it earns no clicks. (The animated variants live in talk-pagestack.tex
% if that ever changes.)
% Requires smtlib.pdf next to the .tex --- the v2.7 reference
% (r2025-07-07); fetch URL in talk-pagestack.tex. PDF page 77 is
% Chapter 5 (Logical Semantics of SMT-LIB Formulas), the part our
% encoding parallels, so it sits in the fan; pdf 55 is a near-blank
% part divider, avoid it.
% 2.3cm: the old 2.0-vs-2.4 ceiling came from the caption line that
% used to sit under the fan; caption gone, the stack earned some of
% the space back (the logic graph beside it gave up width in trade).
\newcommand{\pgstackw}{2.3cm}
\newcommand{\pgstackh}{3.253cm}   % A4 ratio: 2.3 * sqrt(2)
\tikzset{a4page/.style={
	draw=rule, fill=white, line width=0.4pt,
	minimum width=\pgstackw, minimum height=\pgstackh, inner sep=0pt,
	blur shadow={shadow blur steps=10, shadow xshift=1.2pt,
		shadow yshift=-1.6pt, shadow opacity=40}}}
% Two fans, mirror images, centred as a pair with a fixed gutter (they
% originally sat at the ends of the text width bracketing a caption;
% caption cut, gap closed). The left fan pivots on its bottom-left
% corner and opens leftward; the right fan pivots on its bottom-right
% corner and opens rightward, so the two upright top sheets face
% inward. The cover is the left fan's top sheet; every other page is
% just recognisably "dense text". Sheet step is 4 degrees (was 8),
% tightened to pull the deepest sheet's peak corner down --- the peak
% sits at w*sin(step*4) + h*cos(step*4) above the pivot, so most of the
% saving is visual tidiness rather than raw height.
\newcommand{\pagestackfanL}{%
	\begin{tikzpicture}
		\foreach \i/\pg in {5/90, 4/77, 3/45, 2/25} {
			\pgfmathsetmacro{\rot}{(\i-1)*4}
			\node[a4page, rotate around={\rot:(0,0)}, transform shape,
				anchor=south west] at (0,0)
				{\includegraphics[width=\pgstackw, page=\pg]{smtlib.pdf}};
		}
		% Top sheet: the cover ZOOMED, not whole --- drawn entire, its
		% title is ~0.5mm tall. The window is the authors line (the
		% widest band) plus ~20bp of air each side, aspect-matched to
		% the sheet so it fills it exactly, and positioned so the text
		% block (title top to release bottom) is vertically centred in
		% it. Title spans ~75% of the sheet. Trim in bp on the 612x792
		% letter page, measured off a render; re-measure if smtlib.pdf
		% changes revision.
		\node[a4page, anchor=south west] at (0,0)
			{\includegraphics[width=\pgstackw, page=1,
					trim=138bp 173bp 138bp 144bp, clip]{smtlib.pdf}};
	\end{tikzpicture}}
\newcommand{\pagestackfanR}{%
	\begin{tikzpicture}
		\foreach \i/\pg in {5/100, 4/80, 3/60, 2/35, 1/13} {
			\pgfmathsetmacro{\rot}{(\i-1)*4}
			\node[a4page, rotate around={-\rot:(0,0)}, transform shape,
				anchor=south east] at (0,0)
				{\includegraphics[width=\pgstackw, page=\pg]{smtlib.pdf}};
		}
	\end{tikzpicture}}

% ---- the SMT-LIB logic graph ----------------------------------------
% A subset of the standard's logic-inclusion DAG (smt-lib.org draws
% all ~40; the slide needs the shape, not the census). Eight logics:
% the base logics the benchmarks exercise (QF_UF, QF_LIA, QF_AX),
% their joins, and the quantified closures up to AUFLIA. Arrows read
% "is included in". Neutral ink/muted like the pipeline's .smt2 node:
% SMT-LIB is neither the source nor the target language. The lone
% X-crossing in the middle is authentic --- the real graph is a
% thicket, and a planar redraw would misrepresent it politely.
% Sized by the caller: fills \linewidth, so put it in a minipage of
% the width the slide grants it.
\newcommand{\logicgraphfig}{%
	\resizebox{\linewidth}{!}{%
	\begin{tikzpicture}[
		x=1cm, y=1cm,
		logic/.style={font=\footnotesize, text=ink, inner sep=2.5pt},
		ar/.style={-{Stealth[length=4pt]}, draw=muted, thin},
	]
	\node[logic] (qfuf)     at (0, 1.5)     {\code{QF\_UF}};
	\node[logic] (qflia)    at (0, 0)       {\code{QF\_LIA}};
	\node[logic] (qfax)     at (0, -1.5)    {\code{QF\_AX}};
	\node[logic] (qfuflia)  at (2.4, 0.75)  {\code{QF\_UFLIA}};
	\node[logic] (lia)      at (2.4, -0.75) {\code{LIA}};
	\node[logic] (uflia)    at (4.8, 0.75)  {\code{UFLIA}};
	\node[logic] (qfauflia) at (4.8, -0.75) {\code{QF\_AUFLIA}};
	\node[logic] (auflia)   at (7.0, 0)     {\code{AUFLIA}};
	\draw[ar] (qfuf)     -- (qfuflia);
	\draw[ar] (qflia)    -- (qfuflia);
	\draw[ar] (qflia)    -- (lia);
	\draw[ar] (qfax)     to[bend right=10] (qfauflia);
	\draw[ar] (qfuflia)  -- (uflia);
	\draw[ar] (lia)      -- (uflia);
	\draw[ar] (qfuflia)  -- (qfauflia);
	\draw[ar] (uflia)    -- (auflia);
	\draw[ar] (qfauflia) -- (auflia);
	\end{tikzpicture}}}

% ---- proof-tree rule labels (ebproof) --------------------------------
% Quiet, mono, muted: the tree's shape is the point, not the rule names.
\newcommand{\rulelab}[1]{{\color{muted}\scriptsize\code{#1}}}

% ---- the interoperability hub ---------------------------------------
% Ported from talk.tex and restyled: LambdaPi centre in indigo (it is
% the target language), provers in ink, translators as muted mono arrow
% labels, arrows in the pipeline's own style. The PVS arrow was
% unlabelled in the source; left that way rather than guessing a tool.
% 0.88, not 0.92: at 0.92 the hub frame ran 8pt past the text block.
\newcommand{\interopfig}{%
	\resizebox{0.88\textwidth}{!}{%
	\begin{tikzpicture}[
		x=1cm, y=1cm,
		prover/.style={font=\small, text=ink},
		trans/.style={font=\scriptsize, text=muted, fill=white,
				inner sep=1.5pt},
		har/.style={-{Stealth[length=4.5pt]}, draw=muted, semithick},
		hbr/.style={{Stealth[length=4.5pt]}-{Stealth[length=4.5pt]},
				draw=muted, semithick},
	]
	\node[font=\large, text=lpcol] (lp) at (0,0) {\textbf{LambdaPi}};
	\node[prover] (lean) at (-3,2)  {Lean};
	\node[prover] (coq)  at (0,2)   {Rocq};
	\node[prover] (hol)  at (3,2)   {HOL};
	\node[prover] (isa)  at (5,0)   {Isabelle};
	\node[prover] (agda) at (3,-2)  {Agda};
	\node[prover] (pvs)  at (0,-2)  {PVS};
	\node[prover] (k)    at (-3,-2) {$\mathbb{K}$ framework};
	\node[prover] (mat)  at (-5,0)  {Matita};
	\draw[har] (lean) -- (lp) node[trans, midway, sloped, above]
		{\code{lean2dk}};
	\draw[hbr] (coq)  -- (lp) node[trans, midway, sloped, above]
		{\code{vodk}};
	\draw[hbr] (hol)  -- (lp) node[trans, midway, sloped, above]
		{\code{hol2dk}};
	\draw[har] (isa)  -- (lp) node[trans, midway, sloped, above]
		{\code{isabelle\_dk}};
	\draw[har] (agda) -- (lp) node[trans, midway, sloped, above]
		{\code{agda2dk}};
	\draw[har] (pvs)  -- (lp);
	\draw[har] (k)    -- (lp) node[trans, midway, sloped, above]
		{\code{KaMeLo}};
	\draw[hbr] (mat)  -- (lp) node[trans, midway, sloped, above]
		{\code{Krajono}};
	\end{tikzpicture}}}


% Translation-operator notation.
\newcommand{\tra}[1]{\llbracket\,#1\,\rrbracket}
\newcommand{\ty}[1]{{\tra{#1}}_{\mbf{ty}}}
\newcommand{\tm}[1]{{\tra{#1}}_{\mbf{tm}}}

\begin{document}

% Title sits on FLoC cream, like the site's hero band.
\begingroup
\setbeamercolor{background canvas}{bg=floccream}
\frame[plain]{\titlepage}
\endgroup


% =====================================================================
\divider{Part 1 of 3}{Background}
% =====================================================================

% What SMT is. Bullets get the full measure so each is one line; the
% band underneath pairs the solver logos (bullet 1's cast, so always
% shown) with the sledgehammer joke, which lands with its bullet ---
% renumber the \uncover if a bullet is added above it. Logos sit in a
% triangle (cvc5 + z3, veriT below) to keep the band shallow; z3 is
% raised to centre on cvc5's midline, its cap height being taller.
\subsection{Satisfiability modulo theories}

\begin{frame}{\secheader}% was "Satisfiability modulo theories"
	Deciding satisfiability of first-order formulae modulo a theory.
	\begin{itemize}
		\item<+-> \code{cvc5}, \code{veriT}, \code{z3}, and many more\ldots
		\item<+-> ubiquitous in formal methods and automated reasoning.
		\item<+-> backends for theorem provers, e.g.\ Isabelle's
		      \code{sledgehammer}.
	\end{itemize}
	\vfill
	% [c], not [b]: the triangle is shorter than the sledgehammer, so
	% bottom-alignment left their vertical centres misaligned.
	\noindent\makebox[\textwidth]{%
		\begin{minipage}[c]{0.48\textwidth}
			\centering
			\includegraphics[height=8mm]{cvc5-logo}\hspace{7mm}%
			\raisebox{-1.5mm}{\includegraphics[height=11mm]{z3-logo}}\\[3mm]
			\includegraphics[height=8mm]{verit-logo}
		\end{minipage}\hfill
		\begin{minipage}[c]{0.48\textwidth}
			\centering
			\uncover<3->{\includegraphics[height=0.36\textheight]{isabelle-sledgehammer}}
		\end{minipage}}
\end{frame}


\begin{frame}{\secheader}% was "The SMT-LIB standard"
	The \textbf{SMT-LIB standard} gives a common language for solvers.
	\begin{itemize}
		\item<+-> concrete syntax for solver interactions.
		\item<+-> logical semantics based on many-sorted FOL.
		\item<+-> specifications of theories and logics.
	\end{itemize}
	\vfill
	% One stack (the document, top sheet showing the cover zoomed to
	% the text block's width), one logic graph (what it specifies).
	% [c] minipages so the two figures' vertical centres align; the
	% graph fills its minipage, so the 0.51 is what sizes it. The fan
	% column is its natural width and the gutter is fixed --- \hfill
	% plus a too-wide column had put ~23mm between the figures; the
	% makebox centres the pair in the leftover margin.
	\noindent\makebox[\textwidth]{%
		\begin{minipage}[c]{0.29\textwidth}
			\pagestackfanL
		\end{minipage}\hspace{7mm}%
		\begin{minipage}[c]{0.51\textwidth}
			\logicgraphfig
		\end{minipage}}
\end{frame}


\begin{frame}{\secheader}% was "Proof production"
	Proof certificates are not standardized.
	\begin{itemize}
		\item<+-> many competing formats: LFSC, Alethe, DRAT.
		\item<+-> solvers have bespoke inference rules.
	\end{itemize}
	\vfill
	\begin{center}
		\begin{prooftree}
			\hypo{p ∧ ¬ p}
			\infer1[\rulelab{and\_elim}]{p}
			\hypo{p ∧ ¬ p}
			\infer1[\rulelab{and\_elim}]{¬ p}
			\infer2[\rulelab{contra}]{⊥}
		\end{prooftree}
	\end{center}
	\vspace{\gapunit}
\end{frame}


\subsection{Eunoia and Ethos}

\begin{frame}{\secheader}% was "Eunoia and Ethos"
    \eoname{} is a logical framework for SMT proof systems.
	\begin{itemize}
		\item<+-> a signature defines
		      types, constants, programs, rules.
		\item<+-> proofs are checked by C++ tool \code{ethos}.
	\end{itemize}
	\vfill
	% Dictionary entries, per the script's marginal note. IPA and
	% middle dots straight from Fira Sans --- check the log for
	% "Missing character" if these ever change.
	% Stacked cards, 0.85 of the measure and centred: headword + IPA on
	% the first line, definition underneath. Neutral plainbox: Greek
	% words belong to neither language, so the card must not borrow a
	% language tint.
	\noindent\makebox[\textwidth]{%
		\begin{minipage}{0.85\textwidth}
			\begin{plainbox}
				\small
				\textbf{eu·noi·a} \,{\color{muted}/juːˈnɔɪ.ə/}\, \emph{n.}\\
				beautiful thinking; a well-disposed mind.
			\end{plainbox}
			\begin{plainbox}
				\small
				\textbf{e·thos} \,{\color{muted}/ˈiːθɒs/}\, \emph{n.}\\
				character; the guiding spirit of a community.
			\end{plainbox}
		\end{minipage}}
\end{frame}


\begin{frame}[fragile]{\secheader}% was "The Cooperating Proof Calculus"
	The \textbf{co-operating proof calculus} (CPC) is \code{cvc5}'s proof system.
	%
	\begin{itemize}
		% 593 = declare-rule count in cvc5 proofs/eo/cpc excluding
		% expert/ (which adds 27 more), at upstream ed5e0807,
		% 2026-07-15; recount with
		%   grep -rhE '^\s*\(declare-rule' --include='*.eo' \
		%     cpc-upstream-repo/proofs/eo/cpc --exclude-dir=expert | wc -l
		\item<+-> encoded in Eunoia; 593 core rules.
	\end{itemize}
	% The const the rule mentions, then the rule; both verbatim from
	% CPC (theories/Builtin.eo, rules/Booleans.eo), one box, \listsep
	% between fragments as in the paper.
	\begin{eobox}
	\begin{lstlisting}[basicstyle=\lpmono\footnotesize]
(declare-const and (-> Bool Bool Bool)
  :right-assoc-nil true)\end{lstlisting}
	\listsep
	\begin{lstlisting}[basicstyle=\lpmono\footnotesize]
(declare-rule and_elim ((Fs Bool) (i Int))
  :premises (Fs)
  :args (i)
  :conclusion (eo::list_nth and Fs i))\end{lstlisting}
	\end{eobox}
\end{frame}


\begin{frame}{\secheader}% was "LambdaPi"
	\lpname\ implements the $λΠ$-calculus modulo rewriting.
	\gap
	\begin{itemize}
		\item<+-> Successor to Dedukti.
		\item<+-> An interactive proof assistant: tactics, an LSP server.
		\item<+-> A hub for proof interoperability:
	\end{itemize}
	\vfill
	\centerline{\interopfig}
\end{frame}


\begin{frame}{\secheader}% was "Motivation"
	Ethos already checks CPC proofs.
	\gap
	\begin{itemize}
		\item<+-> \textbf{Independence}: a checker outside the
		      \code{cvc5} ecosystem.
		\item<+-> \textbf{Strength}: the calculus itself is
		      typechecked.
		\item<+-> \textbf{Interoperability}: CPC proofs enter the hub.
	\end{itemize}
\end{frame}


% eo2lp closes the on-ramp: the pipeline says what the tool is, the demo
% pair shows it working, and Part 1 ends on the checkmark.
\begin{frame}{\secheader}% was "\code{eo2lp}"
	\begin{center}
		\pipelinefig
	\end{center}
	\begin{itemize}
		\item<+-> Signature translated \textbf{once}; proofs
		      \textbf{per problem}.
		\item<+-> Generic over the signature: no per-rule soundness
		      lemmas, unlike Alethe reconstruction or \textsc{lean-smt}.
		\claim<+-> rule signatures typecheck;
		\claim<+-> programs satisfy subject reduction.
	\end{itemize}
\end{frame}


\begin{frame}[fragile]{\secheader}% was "\code{cvc5} says \code{unsat}"
	\begin{plainbox}
	\begin{lstlisting}[style=smt2, basicstyle=\lpmono\footnotesize]
(set-logic QF_UF)
(declare-const p Bool)
(assert (and p (not p)))
(check-sat)
(get-proof)\end{lstlisting}
	\end{plainbox}
	\onslide<+->
		\eoname\ proof script, from \code{-{}-proof-format=cpc}:
		\begin{eobox}
		\begin{lstlisting}[basicstyle=\lpmono\footnotesize]
(assume @p1 (and p (not p)))
(step @p2 p :rule and_elim :premises (@p1) :args (0))
(step @p3 (not p) :rule and_elim
  :premises (@p1) :args (1))
(step @p4 false :rule contra :premises (@p2 @p3))\end{lstlisting}
		\end{eobox}
\end{frame}


\begin{frame}[fragile]{\secheader}% was "What \code{eo2lp} produces"
	\begin{lpbox}
	\begin{lstlisting}[style=lp, basicstyle=\lpmono\footnotesize]
require open cpc.Cpc;
constant symbol p : τ Bool;
constant symbol {|@p1|} :
  Proof (APP (APP and p)
             (APP (APP and (APP not p)) true));
symbol {|@p2|} ≔ and_elim (of_Z 0) ({|@p1|});
assert ⊢ {|@p2|} : Proof p;
symbol {|@p3|} ≔ and_elim (of_Z 1) ({|@p1|});
assert ⊢ {|@p3|} : Proof (APP not p);
symbol {|@p4|} ≔ contra ({|@p2|}) ({|@p3|});
assert ⊢ {|@p4|} : Proof false;\end{lstlisting}
	\end{lpbox}
	\begin{uncoverenv}<+->
		\yes\ \ \code{lambdapi check}, subject reduction on.
	\end{uncoverenv}
\end{frame}


% =====================================================================
% Part 2 runs Eunoia formalism -> LambdaPi formalism -> translation
% without internal dividers; the box tints (green source, blue target)
% carry the sub-structure.
% =====================================================================
\divider{Part 2 of 3}{Translation}
% =====================================================================

\subsection{Eunoia}

\begin{frame}[fragile]{\secheader}% was "Declarations and attributes"
	One grammar for terms, types and kinds alike.
	\begin{disp}
		\begin{array}{@{}l@{\ }c@{\ }l@{\qquad}l}
			e ∈ \mbf{eo} & ⩴ & s                             & \mcomment{(symbol)}      \\
			             & ∣ & \paren{{s}\ e_1\ \ldots\ e_n} & \mcomment{(application)}
		\end{array}
	\end{disp}
	\onslide<+->
		Attributes determine how $n$-ary applications expand.
		\begin{eobox}
		\begin{lstlisting}[basicstyle=\lpmono\footnotesize]
(declare-const and (-> Bool Bool Bool)
  :right-assoc-nil true)\end{lstlisting}
		\listsep
		\begin{lstlisting}[basicstyle=\lpmono\footnotesize]
(declare-parameterized-const =
  ((A Type :implicit)) (-> A A Bool)
  :chainable and)\end{lstlisting}
		\end{eobox}
\end{frame}


\begin{frame}[fragile]{\secheader}% was "Programs"
	Side-conditions are constants with ordered rewrite cases.
	\begin{eobox}
	\begin{lstlisting}[basicstyle=\lpmono\footnotesize]
(program $to_clause ((F1 Bool) (F2 Bool :list))
  :signature (Bool) Bool
  ( (($to_clause (or F1 F2)) (or F1 F2))
    (($to_clause false)      false)
    (($to_clause F1)         (or F1)) ))\end{lstlisting}
	\end{eobox}
	\begin{uncoverenv}<+->
		Evaluation is by matching, in order, on elaborated arguments.
	\end{uncoverenv}
\end{frame}


\begin{frame}[fragile]{\secheader}% was "Rules"
	A rule's conclusion may be computed by such a program.
	\begin{eobox}
	\begin{lstlisting}[basicstyle=\lpmono\footnotesize]
(declare-rule resolution
  ((C1 Bool) (C2 Bool) (pol Bool) (L Bool))
  :premises (C1 C2) :args (pol L)
  :conclusion ($resolve C1 C2 pol L))\end{lstlisting}
	\end{eobox}
	\begin{uncoverenv}<+->
		Encoding this needs \textbf{computation} in the target.
	\end{uncoverenv}
\end{frame}


\subsection{LambdaPi}

\begin{frame}{\secheader}% was "The $λΠ$-calculus modulo rewriting"
	\begin{itemize}
		\item<+-> Terms:
		      $t ⩴\ x\ ∣\ t_1 ⋅ t_2\ ∣\ \binddot{λ}{x : t_1}{t_2}\ ∣\ \binddot{Π}{x : t_1}{t_2}$
		% Deliberate break: the two forms fill the second line together,
		% rather than the grammar trailing alone after a wrap.
		\item<+-> Symbols, with explicit or implicit parameters:\\
		      $\lpkw{symbol}\ s\ \any{ρ} : t;$
		      \quad $ρ ⩴ (x : t) ∣ [x : t]$
		\item<+-> Computation by rewrite rules:
		      $\lpkw{rule}\ t ↪ t';$
		\item<+-> Typechecking is modulo those rules, and LambdaPi
		      checks they preserve types: \textbf{subject reduction}.
	\end{itemize}
\end{frame}


\begin{frame}[fragile]{\secheader}% was "Types as terms"
	Eunoia types become terms of a universe closed under $(⤳)$.
	\begin{lpbox}
		\begin{lstlisting}[style=lp, basicstyle=\lpmono\footnotesize]
symbol Set : TYPE;          // types as terms
injective symbol τ : Set → TYPE;  // carrier
rule τ ($A ⤳ $B) ↪ τ $A → τ $B;
injective symbol Proof : τ Bool → TYPE;\end{lstlisting}
	\end{lpbox}
	\begin{uncoverenv}<+->
		A Tarski universe: a type-theoretic denotational semantics.
		\par\gap
		\noindent{\small
			\begin{tabular}{@{}l@{\qquad}l@{}}
				\textbf{SMT-LIB models}       & \textbf{$λΠ$ encoding}         \\
				\midrule
				$σ$ denotes a set $⟦σ⟧$       & $T : \lpkw{Set}$               \\
				$⟦f⟧ : ⟦σ_1⟧ × \ldots → ⟦σ'⟧$ & $\tauop\,(T_1 ⤳ \ldots ⤳ T')$ \\
				satisfies the axioms          & typechecks
			\end{tabular}}
	\end{uncoverenv}
\end{frame}


\subsection*{}   % back to the bare section header

\begin{frame}[fragile]{\secheader}% was "Elaboration"
	$\mbf{elab}_{γ} : \mbf{eo} → \mbf{eo}$, with
	$γ$ carrying the attributes declared so far.
	\begin{itemize}
		\item<+-> Default: left-fold into the application symbol
		      $\kw{\_}$.
		      \begin{disp}
			      \mbf{elab}_{γ}\paren{s\ e_1\ \ldots\ e_n}
			      = \paren{\kw{\_}\ \paren{\ldots\paren{\kw{\_}\ s\ e_1}\ldots}\ \ e_n}
		      \end{disp}
		\item<+-> Attributes override it, per the declaration of $s$.
	\end{itemize}
	\begin{eobox}
		\begin{lstlisting}[basicstyle=\lpmono\footnotesize]
(and p (not p))
;; :right-assoc-nil true  =>
(and p (and (not p) true))\end{lstlisting}
	\end{eobox}
	Downstream, every application is binary.
\end{frame}


\begin{frame}[fragile]{\secheader}% was "$\tra{\cdot}_{\mbf{ty}}$ and $\tra{\cdot}_{\mbf{tm}}$"
	Kinds become LambdaPi types; terms and types become terms.
	% Size set outside the environment: \disp opens math mode immediately,
	% and \footnotesize is invalid there (LaTeX warned and half-applied it).
	{\footnotesize
	\begin{disp}
		\begin{array}{@{}r@{\ =\ }l@{\qquad}r@{\ =\ }l}
			\ty{\kw{Type}}                         & \lpkw{Set}
			                                       &
			\tm{e ∗ e'}                            & \tm{e} ⋅ \tm{e'}
			\\[3mm]
			\ty{\paren{\paren{\kw{->}} ∗ e} ∗ e'}  & \ty{e} → \ty{e'}
			                                       &
			\tm{\paren{\kw{->}}}                   & (⤳)
			\\[3mm]
			\ty{e'}                                & \tauop\ \tm{e'}
			                                       &
			\tm{s}                                 & \iden s
		\end{array}
	\end{disp}}
	\onslide<+->
	\exxample A Eunoia type and its image under $\tm{\cdot}$.
	\begin{eobox}
		\begin{lstlisting}[basicstyle=\lpmono\footnotesize]
(-> Bool (BitVec 5))\end{lstlisting}
	\end{eobox}
	\begin{lpbox}
		\begin{lstlisting}[style=lp, basicstyle=\lpmono\footnotesize]
{|Bool|} ⤳ ({|BitVec|} ⋅ {|5|})\end{lstlisting}
	\end{lpbox}
\end{frame}


\begin{frame}[fragile]{\secheader}% was "Declarations become symbols"
	\begin{disp}
		\begin{array}{@{}l}
			\epar{\kw{declare-parameterized-const}\ s\ \epar{\vec ρ}\ e}
			\\[1mm] \quad ⇓ \\[1mm]
			\lpkw{constant symbol}\ {\iden s}\ \tra{ρ_1}\ldots\tra{ρ_n}
			: \tauop\ \tm{e};
		\end{array}
	\end{disp}
	\onslide<+->
		\begin{eobox}
		\begin{lstlisting}[basicstyle=\lpmono\footnotesize]
(declare-parameterized-const =
  ((A Type :implicit)) (-> A A Bool))\end{lstlisting}
		\end{eobox}
		\begin{lpbox}
		\begin{lstlisting}[style=lp, basicstyle=\lpmono\footnotesize]
constant symbol {|=|} [A : Set] :
  τ (A ⤳ A ⤳ Bool);\end{lstlisting}
		\end{lpbox}
	\begin{uncoverenv}<+->
		Implicit parameters become $[\,\cdot\,]$; $γ$ gains
		$\paren{s ↦ α}$.
	\end{uncoverenv}
\end{frame}


\begin{frame}[fragile]{\secheader}% was "Programs become rewrite rules"
	One \lptt{sequential symbol}, one rule per case.
	\begin{lpbox}
	\begin{lstlisting}[style=lp, basicstyle=\lpmono\footnotesize]
sequential symbol {|$from_clause|} :
  τ (Bool ⤳ Bool);
rule {|$from_clause|} (or $F1 $F2)
  ↪ {|eo::ite|} ({|eo::eq|} $F2 false)
      $F1 (or $F1 $F2)
with {|$from_clause|} ($F1) ↪ $F1;\end{lstlisting}
	\end{lpbox}
	\onslide<+->
		A rule becomes an axiom; its conclusion may call the programs.
		\begin{lpbox}
		\begin{lstlisting}[style=lp, basicstyle=\lpmono\footnotesize]
constant symbol and_elim [Fs : τ Bool]
  (i : τ Int) (_ : Proof Fs)
  : Proof ({|eo::list_nth|} and Fs i);\end{lstlisting}
		\end{lpbox}
\end{frame}


\begin{frame}[fragile]{\secheader}% was "The example again"
	\begin{lpbox}
	\begin{lstlisting}[style=lp, basicstyle=\lpmono\footnotesize]
constant symbol {|@p1|} :
  Proof (APP (APP and p)
             (APP (APP and (APP not p)) true));
symbol {|@p2|} ≔ and_elim (of_Z 0) ({|@p1|});
assert ⊢ {|@p2|} : Proof p;\end{lstlisting}
	\end{lpbox}
	\begin{itemize}
		\item<+-> \lptt{true} terminator: from
		      \attrib{right-assoc-nil}.
		\item<+-> \lptt{APP}: from elaboration into $\kw{\_}$.
		\item<+-> The \lptt{assert} holds because
		      $\eo{list\_nth}$ reduces.
	\end{itemize}
\end{frame}


% =====================================================================
\divider{Part 3 of 3}{Results}
% =====================================================================

\begin{frame}{\secheader}% was "\code{eo2lp} and CPC-MINI"
	\begin{itemize}
		\item<+-> ${\sim}5{,}500$ lines of OCaml, built against the
		      LambdaPi API.
		\item<+-> Hand-written \lptt{Prelude.lp}: ${\sim}335$ lines.
		\item<+-> \textbf{CPC-MINI}: 22 modules, ${\sim}5{,}000$ lines,
		      359 rules, 88 programs.
		\item<+-> Logic, arithmetic, arrays, sets, quantifiers, UF.
		\item<+-> Three edits to upstream: pruning, type tightening,
		      bug fixes.
		\item<+-> One-off cost: translated in ${\sim}100$\,ms,
		      checked in ${\sim}160$\,ms.
		\claim<+-> All 22 modules check with subject reduction.
	\end{itemize}
\end{frame}


\begin{frame}{\secheader}% was "Benchmarks"
	1{,}036 \code{cvc5} proofs, 14 fragments, 5\,s timeout.
	\par
	\noindent{\small
		\begin{tabular}{@{}l r r r r r@{}}
			\toprule
			                      & Tot.             & Pass             & \%          & Chk.        & T/O         \\
			\midrule
			Equality              & 200              & 199              & 99          & 0           & 1           \\
			Arrays                & 383              & 383              & 100         & 0           & 0           \\
			Linear int.\ arith.\  & 453              & 420              & 93          & 16          & 17          \\
			\midrule
			\textbf{Total}        & \textbf{1{,}036} & \textbf{1{,}002} & \textbf{97} & \textbf{16} & \textbf{18} \\
			\bottomrule
		\end{tabular}}
	\begin{itemize}
		\item<+-> 0.47\,s per proof to translate and check.
		\item<+-> The 16 check errors are all at the
		      \eott{Real}/\eott{Int} boundary.
		\item<+-> No failure is a semantic disagreement with Ethos.
		\item<+-> Ethos checks in ${<}0.1$\,s and covers all of CPC.
	\end{itemize}
\end{frame}


\begin{frame}{\secheader}% was "What LambdaPi checks that Ethos does not"
	Ethos checks proofs; the encoding also checks the calculus.
	\begin{itemize}
		\item<+-> Every rule signature is well-typed in $T(Σ_{\cpc})$.
		\item<+-> Every program satisfies subject reduction.
	\end{itemize}
	\begin{uncoverenv}<+->
		\gap
		This surfaced bugs in upstream CPC.
		\begin{itemize}
			\item<+-> \eott{\$mk\_cong\_rhs} in \eott{Uf.eo} typed
			      its arguments with the \textbf{codomain} of its
			      function parameter, not the \textbf{domain}.
			\item<+-> Ethos typechecks program clauses only at
			      application time.
		\end{itemize}
	\end{uncoverenv}
\end{frame}


\begin{frame}{\secheader}% was "Conclusion"
	\begin{itemize}
		\item<+-> An automatic translation from \eoname\ to \lpname,
		      at the level of the framework rather than rule by rule.
		\item<+-> 22 modules typechecked; 1{,}002 of 1{,}036 proofs at
		      0.47\,s each.
		\item<+-> Rules become axioms, but LambdaPi typechecks them.
	\end{itemize}
	\begin{uncoverenv}<+->
		\gap
		Next: the rest of CPC; its rules proved sound in $λΠ$;
		SMT-LIB~3.
	\end{uncoverenv}
	\begin{uncoverenv}<+->
		\gap
		\begin{center}
			\lpmono\normalsize github.com/deducteam/eo2lp
		\end{center}
	\end{uncoverenv}
\end{frame}

\end{document}
